% Generated draft: requires mathematical and source review.
\documentclass[11pt]{article}
\usepackage[margin=25mm]{geometry}
\usepackage{fontspec}
\setmainfont{Latin Modern Roman}
\usepackage{amsmath,amssymb,mathtools,unicode-math}
\setmathfont{Latin Modern Math}
\usepackage{enumitem,needspace,xurl,hyperref}
\setlength{\parindent}{0pt}
\setlength{\parskip}{5pt}
\setlength{\emergencystretch}{3em}
\newcommand{\N}{\mathbb{N}}
\newcommand{\Z}{\mathbb{Z}}
\newcommand{\Q}{\mathbb{Q}}
\newcommand{\R}{\mathbb{R}}
\newcommand{\C}{\mathbb{C}}
\newcommand{\F}{\mathbb{F}}
\newcommand{\A}{\mathbb{A}}
\newcommand{\PP}{\mathbb{P}}
\newcommand{\E}{\mathbb{E}}
\newcommand{\eps}{\varepsilon}
\newcommand{\dd}{\,\mathrm{d}}
\begin{document}
\section*{UnsolvedMath ID: 30002870}
\textbf{Source problem code:} catalogue entry
\subsection{Definitions and mathematical statement}
The supplied record concerns a homogeneous space \(X\) and an associated Landau--Ginzburg mirror. A homogeneous space is a space equipped with a transitive group action; the record does not specify the group, stabilizer, ground field, geometric category, or class of homogeneous spaces intended. Write \((Y,W)\) schematically for the mirror space and its potential. This notation does not supply a construction of the mirror or specify whether its data are algebraic or analytic.

Let \(H^\bullet(X)\) denote the cohomology intended by the source, with its theory, coefficients, and grading conventions left unspecified. Write \(H^\bullet_{\mathrm{tw\text{-}dR}}(Y,W)\) for the intended twisted de Rham cohomology. The source does not specify its underlying complex, twisting differential, parameters, or support conditions. It refers to a map
\[
\phi\colon H^\bullet(X)\longrightarrow H^\bullet_{\mathrm{tw\text{-}dR}}(Y,W),
\]
but does not define this map or identify a parameter value at which it is evaluated.

\textbf{Source question.} For a homogeneous space \(X\) and its Landau--Ginzburg mirror in the intended, unspecified class, determine whether:
\begin{enumerate}
\item their relationship admits a more symmetric formulation;
\item the mirror admits coordinates naturally related to the cohomology of \(X\);
\item the indicated cohomological map extends to an isomorphism of bundles with connection.
\end{enumerate}
The first two requests have no supplied mathematical criterion for symmetry, naturality, or the required relation between coordinates and cohomology.

For the third request, the meaning of an isomorphism of bundles with connection can be stated conditionally. Once a common parameter base \(S\), bundles \(\mathcal E_A,\mathcal E_B\) on \(S\), and connections \(\nabla_A,\nabla_B\) have been specified, such an isomorphism is an invertible \(\mathcal O_S\)-linear map
\[
\Phi\colon\mathcal E_A\xrightarrow{\sim}\mathcal E_B
\quad\text{such that}\quad
\nabla_B\circ\Phi=(\Phi\otimes\mathrm{id})\circ\nabla_A.
\]
Here a connection satisfies the Leibniz rule \(\nabla(fs)=f\nabla s+s\otimes\dd f\), with any permitted poles requiring separate specification. Extending \(\phi\) would additionally require specified identifications with fibers or restrictions of these bundles and agreement of \(\Phi\) with \(\phi\) under those identifications. None of these data, nor any boundary or singularity conditions, are supplied. This is therefore a conditional description of the requested compatibility, not a fully specified universal conjecture.
\subsection{Short English statement}
The question asks whether a homogeneous space and its Landau--Ginzburg mirror admit a more symmetric description, coordinates linked naturally to cohomology, and a cohomological correspondence extending to an isomorphism that preserves connections. The intended spaces, cohomology theories, map, parameter base, and connections are unspecified, so the record does not determine a complete precise conjecture. Its partially solved label and claims about selected examples are supplied curation requiring independent review.
\subsection{Sources}
\begin{itemize}
\item Local UnsolvedMath registry: supplied complete source record, permanent ID \(30002870\), source problem code \texttt{catalogue entry}, title \emph{Symmetric Landau--Ginzburg Mirrors of Homogeneous Spaces}; collection: Oberwolfach Reports Open Problems. The project registry locations are \texttt{data/unsolvedmath/problems.json} and the supplementary \texttt{catalog/problems.json}; their contents were not independently inspected for this draft. The record gives proposed year \(2015\) and names no individual proposer.
\item Supplied report citation: \emph{Mirror Symmetry, Hodge Theory and Differential Equations (2016)}, Oberwolfach Reports, DOI \texttt{10.4171/owr/2015/22}, pp.\ 1201--1254. Supplied source URL and literature reference labeled ``Oberwolfach Report 22'': \url{https://doi.org/10.4171/owr/2015/22}. The record describes this collection as documenting workshop discussions and open problems. The report citation was supplied and was not independently consulted.
\item Supplied literature assessment, checked \textbf{2026-08-21}: status \texttt{partially\_solved}. The assessment states that the mirror-symmetry and A/B-model connection is proved for selected homogeneous examples such as Grassmannians and quadrics, while the general conjectural bundle-with-connection statement remains open. This is supplied dated curation, not an independently verified mathematical or literature assessment. Its listed literature source is \url{https://doi.org/10.4171/owr/2015/22}; the supplied background expressly requests independent verification of time-sensitive status and cited claims.
\item The registry labels the statement \texttt{corrected\_verified} and reports a manual semantic, typographic, and LaTeX review of its cleanup. These are supplied editorial qualifications and do not establish the missing mathematical data or constitute independent verification here.
\end{itemize}
\end{document}
